\section{SF-E829. A wedge-pairing gate for arbitrary surface
directions}\label{SF-E829:sf-e829.-a-wedge-pairing-gate-for-arbitrary-surface-directions}

Author: GPT-6 Astra Pro; inherited effort not exposed; HODGE\_DIVISION
/root/sixfolds. Date: 2026-11-21.

\textbf{Status:} independently whole-reviewed by /root and
/root/g34\_eightfolds; both mathematical reviews PASS on the frozen draft
\texttt{4cab77bfcc30b75151e946b3a1da06e29e2461ce59f8787c326c9aa5ab666b28}.
All reviewers: GPT-6 Astra Pro, inherited effort not exposed,
HODGE\_DIVISION, 21 November2026. The peer also supplied an independently
written supplementary exact exterior certificate; the infinite theorem is
the complete proof below, not an extrapolation from finite checks. This
extends E827 beyond its rational four-divisor algebra. It is a necessary
support-incidence gate for the same actual equivariant extension
architecture. It is not a classification of all coherent sheaves, an
existence theorem for a semiregular seed, or an assertion that the remaining
elliptic/surface configurations succeed.

\subsection{1. Exact statement}\label{SF-E829:exact-statement}

Let X be a polarized abelian fourfold and retain the actual coherent
extension and actual equivariant blocks from E827:

\[
 0\longrightarrow A\longrightarrow E\longrightarrow B\longrightarrow0,
 \qquad A=A_0\oplus\bigoplus_{a=1}^u L_a^{n_a},
 \qquad B=\bigoplus_{b=1}^v M_b^{m_b}.
 \tag{E829.1}
\]

Every cross difference line N\_ab=L\_a tensor M\_b inverse is assumed to
have nonzero first Chern class. This is automatic for the opposite slopes in
E826. Write rank(N) for the complex rank of its Hermitian form, relative to
the fixed positive polarization. For every pair, require either rank(N\_ab)
at least2 or \(H^*(X,N_{ab})=H^*(X,N_{ab}^{-1})=0\). Thus a rank-one pair is
allowed only if its flat kernel twist makes it cohomologically irrelevant.
The rank-one cases with nonzero cohomology are deliberately not covered by
the theorem.

Put V=H1(X,O\_X), dim V=4. For each rank-two difference whose nondegenerate
surface part has index one and whose flat kernel part is trivial after the
product isogeny below, let V\_N be the two-dimensional subspace of V coming
from that surface factor. Let z\_N=Lambda2 V\_N, a line in Lambda2 V. Define

\[
 Z=\operatorname{span}_\mathbf C\{z_N\}\subset\Lambda^2V,
 \qquad z=\dim Z.
 \tag{E829.2}
\]

The construction below gives these subspaces canonically after choosing the
fixed polarization. Including an extra acyclic rank-two pair in Z is
permitted and only weakens the bound. Choose a nonzero volume in Lambda4 V
and let q be the nondegenerate symmetric wedge pairing on Lambda2 V. Let rho
be the dimension of the radical of q restricted to Z. Put

\[
 w=6-z-\rho.
 \tag{E829.3}
\]

Let J be any subset of negative-block indices for which Ext3(M\_b,A0)=0. For
every specified coherent effective finite action preserving the actual
sequence and blocks,

\[
 \operatorname{ext}^2_\Gamma(E,E)\ge w(u+|J|).
 \tag{E829.4}
\]

In particular, if the complete normalized character is a nonzero multiple of
the balanced v\_q of E801, with any rational algebraic B field, weak
semiregularity requires w(u+\textbar J\textbar) at most18. The letter q in
the wedge pairing and the parameter in v\_q denote different objects; in the
proof below the pairing is written q\_wedge when needed to avoid ambiguity.

\subsection{2. The product isogeny and the cohomology
calculation}\label{SF-E829:the-product-isogeny-and-the-cohomology-calculation}

We give the needed degenerate-line reduction directly, including the isogeny
and its scope. On the real tangent space of X, the alternating form E\_theta
of the polarization identifies the alternating form E\_N of a line bundle
with a rational linear operator T:

\[
 E_N(x,y)=E_\theta(Tx,y).
 \tag{E829.5}
\]

Both forms are integral on the lattice, and E\_theta is nondegenerate, so T
is rational. Their type(1,1) condition implies that T commutes with the
complex structure. Alternation gives E\_theta(Tx,y)=E\_theta(x,Ty), so T is
self-adjoint for the positive Hermitian form associated with theta. It is
therefore diagonalizable with real eigenvalues. Its kernel K\_R and image
C\_R are rational complex subspaces, and are E\_theta-orthogonal
complements. Intersecting them with the lattice produces complex subtori K
and C. The restriction of theta to each is positive, so these are abelian
subvarieties. Their sum map

\[
 f:C\times K\longrightarrow X
 \tag{E829.6}
\]

is an isogeny: the differential is an isomorphism, and the sum of the two
intersection lattices has finite index in the original lattice. This is a
lattice proof of the needed complement; no generic splitting or rational RM
eigenspace is being treated as an abelian subvariety.

The complex dimension of C is rank(N), and E\_N is nondegenerate on C and
zero on K and on the mixed terms. Thus

\[
 f^*N=N_C\boxtimes P_K,
 \tag{E829.7}
\]

where N\_C is a nondegenerate line bundle and P\_K is a degree-zero line
bundle. To see the line-level assertion, subtract the pullback of the
restriction to C: the remaining line has zero first Chern class, and a flat
character of the product lattice is a product of characters of its two
factors. Absorb its C component into N\_C. This proves(E829.7), not merely
its numerical counterpart.

A nontrivial degree-zero line on an abelian variety has no cohomology. One
may use the previously read published Pic0 vanishing, or the direct
unitary-flat proof: the corresponding flat character is nontrivial, and the
flat Laplacian on its torus has no zero-frequency harmonic mode. Hence if
P\_K is nontrivial, both lines in(E829.7) and their inverses are acyclic. If
P\_K is trivial, the abelian index theorem puts the cohomology of N\_C in
exactly one degree j, and the inverse in degree r-j, where r=dim C. Kunneth
then describes all cohomology and all its products.

The pullback of H1(O\_X) is an isomorphism onto H1(O\_C) direct sum
H1(O\_K), so it identifies V with these two summands. Define V\_N=H1(O\_C)
under this identification. It is independent of lattice indices in the above
construction because C\_R=im T is fixed. The pullback of \(H^*(X,N)\) is
injective by finite-etale trace: trace after pullback is multiplication by
deg f.~Products commute with pullback. Thus every vanishing and every
product-containment proved on C times K descends to X. This use of an
isogeny concerns cohomology and products only; no action, object, or
semiregularity is being pulled back in the conclusion.

\subsection{3. Exhaustive ranks at least
two}\label{SF-E829:exhaustive-ranks-at-least-two}

Let V\_C=H1(O\_C), V\_K=H1(O\_K), with dimensions r and k=4-r. When P\_K is
trivial, the only possible degree-one terms are

\[
 H^1(f^*N)=
 \begin{cases}
 H^0(N_C)\otimes V_K,&j=0,\\
 H^1(N_C)\otimes\mathbf C,&j=1,\\
 0,&j\ge2.
 \end{cases}
 \tag{E829.8}
\]

Multiplication by any class containing a V\_C factor raises the C
cohomological degree away from its sole index and is therefore zero. Only
the component in Lambda2 V\_K can act nontrivially from degree1 to 3. This
observation also applies to N inverse, at its index r-j.

If r=4, the line is nondegenerate on X. Its degree-one group and that of its
inverse cannot both be nonzero, and its own degree-three group is zero
whenever its degree-one group is nonzero. Thus all diagonal degree-one
products vanish and every element of Lambda2 V kills both degree-one groups
into degree3.

If r=3, then k=1 and Lambda2 V\_K=0. Every element again kills degree1 into
3. The indices of N\_C and its inverse sum to3, so both cannot have index at
most1. Thus both degree-one groups cannot occur simultaneously, and all
their diagonal products vanish. This includes the j=0 term with its
one-dimensional kernel factor; no nondegeneracy on the whole X is assumed.

If r=2 and j=0 or2, only one of the two lines can have degree-one
cohomology. The possible group has j=0 and is H0(N\_C) tensor V\_K.
Multiplication by Lambda2 V\_K lands in Lambda3 V\_K=0 because k=2. The
C-containing components already vanish. Thus all diagonal products vanish
and every element of Lambda2 V kills degree1 into3 in these cases.

The sole remaining case with r at least2 is r=2,j=1. Here both lines have

\[
 H^1(f^*N)=H^1(N_C)\otimes\mathbf C,
 \qquad
 H^3(f^*N)=H^1(N_C)\otimes\Lambda^2V_K.
 \tag{E829.9}
\]

Their diagonal degree-one Yoneda products lie in Lambda2 V\_C, the support
line z\_N. Multiplication of a degree-one class by w0 in Lambda2 V depends
only on its projection to Lambda2 V\_K. It is zero whenever that projection
vanishes. Since both V\_C and V\_K have dimension2, this condition is
exactly

\[
 w_0\wedge z_N=0,
 \tag{E829.10}
\]

meaning the wedge with any generator of that line is zero. Equivalently, w0
is orthogonal to z\_N for q\_wedge. This completes the actual cup-product
analysis of every pair allowed in the theorem.

\subsection{4. The surviving diagonal space and invariant exact
sequences}\label{SF-E829:the-surviving-diagonal-space-and-invariant-exact-sequences}

For all cross pairs together, the preceding calculations show that every
incoming diagonal degree-one product lies in Z, and every element of Z
perpendicular kills both relevant degree-one groups into degree3. The
pairing on Lambda2 V is nondegenerate. Thus

\[
 \dim Z^\perp=6-z,
 \qquad Z\cap Z^\perp=\operatorname{rad}(q_{\mathrm{wedge}}|_Z).
 \tag{E829.11}
\]

Choose any vector-space complement W to this radical inside Z perpendicular.
Then W intersects Z trivially, and dim W=6-z-rho=w. All such subspaces are
fixed by the allowed translation/tensor functors on H2(O\_X), since those
functors act trivially on this cohomology. No invariant choice beyond this
trivial action is needed.

Now apply exactly the actual diagonal-survival proof E827.1--5: the
A-diagonal W identity classes inject through Hom(A,triangle); their next
boundary vanishes because W kills Ext1 into Ext3. The selected B-diagonal
classes lift because Ext3(M\_b,A0)=0. Normalized diagonal matrix trace
modulo Z kills the preceding boundary, with its A0 term zero by Serre
duality. Finite invariants are exact and all selected classes are invariant.
Consequently their contributions add to w(u+\textbar J\textbar). This
is(E829.4). There is no degeneration claim for a filtered Ext spectral
sequence and no assumption that different Chern-character terms are summands
of E.

For the full balanced character, its cap rank is18 and the rational B shear
is invertible. The already reviewed full supplied S21 adjunction therefore
forces invariant Ext2 dimension18 under weak semiregularity.
Equation(E829.4) gives the claimed necessary inequality.

\subsection{5. Concrete geometric consequences and the elliptic
limitation}\label{SF-E829:concrete-geometric-consequences-and-the-elliptic-limitation}

With no exceptional rank-two index-one pair, Z=0 and w=6. Thus E827's
six-per-block obstruction extends to every rank-three degeneration and to
every rank-two index-zero or index-two degeneration, without assuming the
four-divisor algebra. For u at least4 the weak balanced seed is excluded
under these hypotheses, even if the unit block has nonzero opposite Ext3.

With a single exceptional support line, z=1 and rho=1 because a decomposable
two-vector squares to zero. Hence w=4. With exactly two independent support
lines generated by z1,z2, there are two alternatives:

\begin{itemize}
\tightlist
\item
  if z1 wedge z2 is nonzero, the restriction pairing has matrix
  {[}{[}0,c{]},{[}c,0{]}{]}, c nonzero, so z=2,rho=0 and w=4;
\item
  if z1 wedge z2 is zero, the restriction is zero, so z=2,rho=2 and w=2.
\end{itemize}

For two-dimensional subspaces of a four-dimensional V, the second case means
the two support planes intersect in a nonzero vector; the first means they
are complementary. These are the actual rational surface support directions
defined from their numerical line forms, not the irrational RM embedding
planes. In particular the two complementary surface directions of E827's
four-divisor algebra give w=4, recovering that earlier bound. Arbitrary
additional directions are controlled by the explicit formula6-dimZ-dimrad,
not by an unstated general-position claim.

The rank-one limitation is genuine for this method. On a product C times K
with C elliptic and K three-dimensional, take a nontrivial positive line
N\_C and trivial P\_K. Then N\_C has index0 and its inverse index1. The two
degree-one spaces are H0(N\_C) tensor V\_K and H1(N\_C inverse) tensor C. By
elliptic Serre duality their diagonal products span

\[
 Z_{\mathrm{ell}}=V_C\wedge V_K.
 \tag{E829.12}
\]

For multiplication to kill the inverse degree-one group into degree3, a
two-vector must have zero Lambda2 V\_K component. Since V\_C is
one-dimensional, the entire allowable annihilator is exactly the same space
V\_C wedge V\_K. It has zero possible complement disjoint from Z\_ell. Thus
the E827 lemma cannot extract any nonzero common W from this single pair.
This is a demonstrated limitation of the diagonal method, not a proof of
semiregularity or existence of a balanced seed. If a rank-one pair is
acyclic because of a nontrivial flat kernel factor, all its products vanish
and it causes no restriction, as permitted in the statement.

\subsection{6. Sources and check boundary}\label{SF-E829:sources-and-check-boundary}

The product isogeny and its numerical/line-level decomposition are proved
above from the rational self-adjoint operator of the two Riemann forms. The
only substantive external cohomology input is the original Mumford Abelian
Varieties(1970), section16 printedpp150--153 and162--163, already fully read
in the invoked scope for E826, together with the previously read Pic0
vanishing or the direct flat proof supplied above. Kunneth, finite trace and
the actual cup products are stated with their full uses. The actual
triangle/invariant argument is the wholly proved E827 lemma; the balanced
rank consequence retains the exact supplied S21 and published BF status
already recorded there. No new numerical experiment, period rank, finite
action or semiregular object is asserted.
